\section{Simulation data, samples, and observables}
\label{sec:data}

\subsection{TNG300 and the galaxy sample}

The analysis uses the TNG300 hydrodynamical simulation from the IllustrisTNG
suite \citep{Pillepich2018,Nelson2019}. TNG300 provides a large volume together
with baryonic galaxy formation physics, halo catalogs, and merger trees. The
present results are therefore conditional on one simulation's numerical
resolution and subgrid model; no claim of transfer to another simulation or to
observations is made in these experiments.

The objects are massive central galaxies and their main-branch halos. The
samples used by each stage are summarized in \cref{tab:samples}. The reduction
from the single-epoch to the five-epoch sample is caused by the requirement of
valid progenitor profiles at every fitted snapshot and the removal of two known
broken progenitor matches. Concentration analyses use a further intersection
with valid fitted concentration histories. The one-object difference between
the \ExpFifty{} and \ExpFiftyOne{} concentration samples follows from applying
the five-epoch progenitor quality requirement in \ExpFiftyOne{}.

\begin{table}[htbp]
  \centering
  \caption{Samples used in Experiments 50 and 51. A valid CoG is a finite,
  monotone cumulative stellar mass profile on the adopted radial grid.}
  \label{tab:samples}
  \begin{tabularx}{\textwidth}{@{}lrlX@{}}
    \toprule
    Analysis & $N$ & Redshift(s) & Additional requirement \\
    \midrule
    \ExpFifty{} main sample & 2539 & 0.4 & Valid \diffmah{}, $\ctwo$, and CoG \\
    \ExpFifty{} concentration history & 2366 & 0.4--2.0 & Five valid concentration measurements \\
    \ExpFiftyOne{} main sample & 2395 & 0.4--2.0 & Five valid progenitor CoGs; broken matches removed \\
    \ExpFiftyOne{} epoch-local sample & 2365 & 0.4--2.0 & Valid truncated \diffmah{} and concurrent $\ctwo$ at all epochs \\
    \bottomrule
  \end{tabularx}
\end{table}

The five profile epochs are $z=(0.4,0.7,1.0,1.5,2.0)$. All held-out model
comparisons use five shared galaxy folds. Sharing folds matters because the
difference between two competitive models is much smaller than their absolute
prediction error.

\subsection{Curve of growth}

The primary observable is the projected, CoG-derived cumulative stellar mass
\begin{equation}
  M_\star(<R)=2\pi\int_0^R \Sigma_\star(R')R'\,dR',
  \label{eq:cog_definition}
\end{equation}
sampled at 24 radii from $2.0$ to approximately $148.2\kpc$. Its normalized
form is
\begin{equation}
  F_\star(R)=\frac{M_\star(<R)}{M_{\star,\rm tot}},\qquad
  M_{\star,\rm tot}\equiv M_\star(<148.2\kpc).
  \label{eq:normalized_cog}
\end{equation}
The measured CoG is monotone by construction. Modeling the cumulative profile
is useful because central apertures, radial annuli, total mass, and enclosed-mass
radii can all be derived from the same prediction. The direct projected-particle
aperture masses are reserved for cross-checks; they are not substituted for the
CoG-derived observable.

Each galaxy profile is a single projected realization of a generally triaxial
and sometimes disturbed system. Projection, merger geometry, and baryonic
history therefore contribute irreducible scatter relative to the selected halo
features.

\subsection{Halo-side measurements and units}

The mass history is the main-branch peak halo mass, $\Mpeak(t)$. Masses used by
the emulator are $h$-free and reported as $\logten(M/\Msun)$. The official
\diffmah{} comparison catalog stores mass in $\Msun/h$; on ingestion the analysis
adds
\begin{equation}
  \Delta\logten M=\logten(1/h)=0.169\dex
\end{equation}
for the simulation value of $h$. Omitting this conversion would produce a
numerically plausible but physically spurious offset comparable to several of
the effects studied here.

The secondary halo property is the spherical-overdensity concentration
\begin{equation}
  \ctwo=\frac{\Rtwo}{r_s},
  \label{eq:concentration}
\end{equation}
where $\Rtwo$ encloses a mean density 200 times the critical density and $r_s$
is the scale radius of an NFW-like profile \citep{Navarro1997}. The epoch at
which $\ctwo$ is measured is part of the feature definition.

\subsection{Information boundary at high redshift}

A high-redshift prediction can be posed in two ways. If the objects are selected
as descendants at $z=0.4$, their full later halo histories are known and may be
used to reconstruct their progenitor profiles. This is a legitimate conditional
history model. It is not, however, a portable model for an independently
selected $z=2$ halo catalog. For that use, every halo feature must be measurable
using only times $t\leq t(z)$. This distinction becomes a central experimental
control in \ExpFiftyOne{}.
